\documentclass[11pt,letterpaper]{article}

\usepackage[top=1in, bottom=1in, left=1in, right=1in, headheight=14pt]{geometry}
\usepackage{fontspec}
\setmainfont{DejaVu Serif}
\setsansfont{DejaVu Sans}
\setmonofont{DejaVu Sans Mono}

\usepackage{xcolor}
\usepackage{titlesec}
\usepackage{fancyhdr}
\usepackage{enumitem}
\usepackage{hyperref}
\usepackage{booktabs}
\usepackage{tabularx}
\usepackage{longtable}
\usepackage{colortbl}
\usepackage{multirow}
\usepackage{array}
\usepackage{parskip}
\usepackage{tcolorbox}
\usepackage{graphicx}
\usepackage{lastpage}
\usepackage{pdfpages}

% === COLOR SCHEME: PNW Ecology (Forest/River/Earth) ===
\definecolor{primary}{HTML}{1B5E20}
\definecolor{secondary}{HTML}{1565C0}
\definecolor{tertiary}{HTML}{4E342E}
\definecolor{accent}{HTML}{2E7D32}
\definecolor{lightprimary}{HTML}{E8F5E9}
\definecolor{lightsecondary}{HTML}{E3F2FD}
\definecolor{lighttertiary}{HTML}{EFEBE9}
\definecolor{rowgray}{HTML}{F5F5F5}
\definecolor{bordergray}{HTML}{BDBDBD}
\definecolor{subtlegray}{HTML}{757575}
\definecolor{darktext}{HTML}{212121}

% === SECTION FORMATTING ===
\titleformat{\section}
  {\Large\bfseries\sffamily\color{primary}}
  {\thesection}{1em}{}
  [\vspace{-0.5em}\textcolor{bordergray}{\rule{\textwidth}{0.4pt}}]

\titleformat{\subsection}
  {\large\bfseries\sffamily\color{secondary}}
  {\thesubsection}{1em}{}

\titleformat{\subsubsection}
  {\normalsize\bfseries\sffamily\color{tertiary}}
  {\thesubsubsection}{1em}{}

\titlespacing*{\section}{0pt}{1.5em}{0.8em}
\titlespacing*{\subsection}{0pt}{1.2em}{0.5em}
\titlespacing*{\subsubsection}{0pt}{0.8em}{0.3em}

% === CUSTOM ENVIRONMENTS ===
\newcommand{\stageheader}[2]{%
  \noindent\colorbox{#2}{\makebox[\textwidth][l]{%
    \textcolor{white}{\bfseries\sffamily\hspace{6pt}#1\hspace{6pt}}}}%
  \vspace{0.5em}\par%
}

\newtcolorbox{asciibox}{
  colback=rowgray, colframe=bordergray,
  arc=1pt, boxrule=0.5pt,
  left=6pt, right=6pt, top=4pt, bottom=4pt,
  fontupper=\small\ttfamily,
  breakable
}

\newtcolorbox{quotebox}{
  colback=lightprimary, colframe=primary,
  arc=2pt, boxrule=0.5pt,
  left=12pt, right=12pt, top=8pt, bottom=8pt,
  breakable
}

\newcommand{\metafield}[2]{\textbf{\sffamily #1} #2\par}
\newcolumntype{L}[1]{>{\raggedright\arraybackslash}p{#1}}
\newcolumntype{C}[1]{>{\centering\arraybackslash}p{#1}}
\newcommand{\tableheader}[1]{\textbf{\sffamily\textcolor{white}{#1}}}

% === HEADERS/FOOTERS ===
\pagestyle{fancy}
\fancyhf{}
\fancyhead[L]{\small\sffamily\textcolor{subtlegray}{PNW Ecosystem Research}}
\fancyhead[R]{\small\sffamily\textcolor{subtlegray}{GSD Mission Package}}
\fancyfoot[C]{\small\sffamily\textcolor{subtlegray}{Page \thepage\ of \pageref{LastPage}}}
\renewcommand{\headrulewidth}{0.4pt}
\renewcommand{\footrulewidth}{0pt}

% === HYPERLINKS ===
\hypersetup{
  colorlinks=true,
  linkcolor=secondary,
  urlcolor=secondary,
  pdfauthor={GSD Ecosystem / Tibsfox},
  pdftitle={PNW Living Systems -- Vision to Mission Pack},
  pdfsubject={Vision to Mission Pipeline},
}

\begin{document}

% ============================================================
% TITLE PAGE
% ============================================================
\thispagestyle{empty}
\begin{center}
\vspace*{2cm}

{\small\sffamily\textcolor{primary}{\textbf{GSD RESEARCH MISSION PACKAGE}}}

\vspace{0.3cm}
\textcolor{bordergray}{\rule{8cm}{0.4pt}}

\vspace{1.2cm}
{\fontsize{26}{32}\selectfont\bfseries\sffamily\textcolor{primary}{Pacific Northwest\\[0.3em]Living Systems\\[0.3em]Taxonomy \& Chipset}}

\vspace{0.6cm}
{\Large\sffamily\textcolor{secondary}{From Tahoma's Summit to the Salish Sea Floor}}

\vspace{0.4cm}
{\large\sffamily\itshape\textcolor{tertiary}{A Deep Research \& Engineering Study}}

\vspace{0.4cm}
{\normalsize\sffamily\textcolor{subtlegray}{Math Co-Processor $\rightarrow$ Taxonomy Map-Reduce $\rightarrow$ Chipset Derivation $\rightarrow$ Engineering Pass}}

\vspace{2cm}
{\sffamily\textcolor{accent}{Vision $\rightarrow$ Research $\rightarrow$ Mission}}\\[0.3em]
{\small\sffamily\textcolor{accent}{Full Pipeline Execution with Minecraft Simulation Layer}}

\vspace{2cm}
{\small\sffamily\textcolor{subtlegray}{Date: March 7, 2026~~|~~Status: Ready for GSD Orchestrator}}\\[0.3em]
{\small\sffamily\textcolor{subtlegray}{Activation Profile: Fleet~~|~~Estimated: 42 context windows across 14 sessions}}\\[0.3em]
{\small\sffamily\textcolor{subtlegray}{Integrates: Heritage Skills, Cooking with Claude, College Structure, GSD-OS}}
\end{center}
\newpage

% ============================================================
% TABLE OF CONTENTS
% ============================================================
\tableofcontents
\newpage

% ============================================================
% STAGE 1 --- VISION DOCUMENT
% ============================================================
\stageheader{STAGE 1 --- VISION DOCUMENT}{primary}

\section{Pacific Northwest Living Systems --- Vision Guide}

\metafield{Date:}{March 7, 2026}
\metafield{Status:}{Initial Vision / Pre-Research}
\metafield{Depends On:}{gsd-chipset-architecture-vision.md, gsd-silicon-layer-spec.md, gsd-skill-creator-analysis.md, 04-salish-sea-expansion-vision.md, cooking-with-claude-compiled-session.md}
\metafield{Context:}{Comprehensive PNW ecosystem research organized through GPU-accelerated taxonomy processing, derived into a GSD chipset architecture, integrated with the College Structure and Heritage Skills educational pack, and rendered as a Minecraft simulation world scaled to geographic accuracy.}

\subsection{Vision}

Stand on the summit of Tahoma --- Mount Rainier --- at 14,410 feet and look in every direction. To the west, the land descends through subalpine meadows, old-growth Douglas fir forests, and river valleys to the shores of Puget Sound, whose glacially carved basins plunge 930 feet below the surface. To the north, the Cascade Range stretches toward the Canadian border, its western slopes draining through salmon rivers into the Salish Sea. To the south, the volcanic arc continues through Mount St.\ Helens, Mount Adams, and Mount Hood. To the east, rain shadow transforms wet forest into dry steppe within miles.

This vertical transect --- 15,340 feet from the deepest point of Puget Sound to the summit ice cap --- contains one of the most biologically dense and ecologically complex gradients on Earth. Every 1,000 feet of elevation change brings shifts in plant communities, animal assemblages, fungal networks, and microclimates equivalent to hundreds of miles of latitudinal change. The salmon that spawn in the headwater streams of Rainier's glacial rivers carry marine-derived nutrients from the Pacific Ocean into the alpine forests, feeding the trees that stabilize the slopes that filter the water that sustains the salmon. The western red cedar that provides the material culture of the Coast Salish peoples grows in the same watersheds where the orcas hunt the same salmon that the Lummi reef nets harvested for millennia.

This pack applies the GSD ecosystem's architectural principles to the study of these living systems. It uses the GPU-based math co-processor (from the Silicon Layer specification) to accelerate taxonomy processing. It organizes research through map-reduce parallel execution. It derives a chipset architecture from the taxonomy itself --- each major taxonomic group becoming a specialized chip, each ecological relationship becoming a bus route. It integrates the results with the College Structure's biology department and the Heritage Skills educational pack's PNW expansion. And it renders the entire system as a Minecraft world scaled so that Tahoma's summit sits at build height (y=319) and the deepest point of Puget Sound sits at bedrock (y=-64).

The Amiga Principle applies: architectural leverage over raw computational power. We do not brute-force 10,000 species descriptions sequentially. We organize the taxonomy so that shared research --- habitat types, elevation bands, ecological roles --- is computed once and reused across every species that shares those attributes. The math co-processor handles the geometric calculations (elevation gradients, watershed topology, population dynamics). The chipset routes research queries to the right specialist agent. The result is a comprehensive PNW ecosystem research block built in a fraction of the time brute-force would require.

\subsection{Problem Statement}

\begin{enumerate}[leftmargin=2em]
\item \textbf{Research duplication across species.} A naive approach to PNW ecosystem documentation would research each species independently, duplicating habitat descriptions, elevation band characteristics, and ecological context thousands of times. This wastes tokens and produces inconsistent results. An organized taxonomy with shared attribute layers eliminates this redundancy.

\item \textbf{No computational acceleration for ecological modeling.} Current GSD research missions treat the GPU as a future concern. The PNW ecosystem's spatial and mathematical complexity --- watershed topology, elevation gradients, population dynamics, nutrient cycling --- demands the math co-processor described in the Silicon Layer spec. This mission activates it.

\item \textbf{Disconnection between ecological knowledge and chipset architecture.} The GSD chipset architecture uses hardware metaphors abstractly. This mission makes them concrete: each taxonomic domain becomes a literal chip, each ecological relationship becomes a bus, each research query becomes a routed instruction. The taxonomy \emph{is} the chipset.

\item \textbf{No spatial simulation layer.} The Heritage Skills pack teaches PNW ecology through text and badges. A Minecraft world scaled to geographic accuracy --- where you can walk from the Puget Sound floor to Rainier's summit and encounter the correct biomes at the correct elevations --- provides embodied learning that text cannot.

\item \textbf{Fragmented college departments.} The College Structure has a biology wing conceptually, but no concrete seed data. This mission produces the foundational dataset: a complete PNW taxonomy organized as explorable code, cross-referenced with the math department (population dynamics, fractal branching in river systems) and the heritage skills department (ethnobotany, traditional ecological knowledge).
\end{enumerate}

\subsection{Core Concept}

\begin{quotebox}
\textbf{One taxonomy, four passes, one chipset.} Research the PNW ecosystem through a staged pipeline: (1) initialize the math co-processor for spatial/ecological computation, (2) execute a full taxonomy map-reduce producing comprehensive species profiles with shared attribute layers, (3) derive a chipset architecture where taxonomic groups become chips and ecological relationships become buses, (4) run an engineering optimization pass to minimize redundancy and maximize throughput. The output is both a research document and a functioning GSD chipset definition.
\end{quotebox}

\subsubsection{The Four-Pass Pipeline}

\begin{asciibox}
PASS 1: MATH CO-PROCESSOR INITIALIZATION
=========================================
  GPU Math Engine (Silicon Layer)
    |-- Elevation Gradient Calculator (14,410 ft summit to -930 ft seabed)
    |-- Watershed Topology Engine (river network graph)
    |-- Population Dynamics Solver (predator-prey, carrying capacity)
    |-- Fractal Geometry (L-systems for branching: rivers, trees, fungi)
    |-- Spatial Projection (lat/long to Minecraft block coordinates)
    |
    v
PASS 2: TAXONOMY MAP-REDUCE
============================
  SHARED ATTRIBUTE LAYERS (computed once, reused everywhere)
    |-- Elevation Bands (Alpine > Subalpine > Montane > Lowland > Marine)
    |-- Habitat Types (Old-growth, Riparian, Estuary, Pelagic, Intertidal...)
    |-- Ecological Roles (Keystone, Apex Predator, Decomposer, Pollinator...)
    |-- Conservation Status (ESA listing, state status, tribal significance)
    |-- Cultural Significance (Heritage Skills cross-reference)
    |
    +-- MAP: Research agents fan out across taxonomy branches
    |     |-- Flora Agent --> Conifers, Deciduous, Understory, Marine Plants
    |     |-- Fauna Agent --> Mammals, Birds, Fish, Amphibians, Reptiles
    |     |-- Fungi Agent --> Mycorrhizal Networks, Decomposers, Lichens
    |     |-- Marine Agent --> Invertebrates, Algae, Plankton, Marine Mammals
    |     |-- Microbiome Agent --> Soil, Aquatic, Symbiotic Communities
    |
    +-- REDUCE: Synthesis agents merge by ecological relationship
          |-- Predator-Prey Networks
          |-- Nutrient Cycling Pathways
          |-- Symbiotic Relationships
          |-- Watershed Connectivity
          |-- Cultural-Ecological Intersections
    |
    v
PASS 3: CHIPSET DERIVATION
===========================
  Taxonomy --> Chipset Architecture
    |-- Each major group = a CHIP (Flora-Chip, Fauna-Chip, etc.)
    |-- Each ecological role = an INSTRUCTION TYPE
    |-- Each habitat type = a MEMORY BANK
    |-- Each relationship = a BUS ROUTE
    |-- Each elevation band = a CLOCK DOMAIN
    |
    v
PASS 4: ENGINEERING OPTIMIZATION
=================================
  Optimize derived chipset for:
    |-- Cache coherence (shared habitat data across chips)
    |-- Bus contention (minimize cross-chip communication)
    |-- Power budget (token allocation across research agents)
    |-- Throughput (parallel execution across elevation bands)
    |-- Integration (College dept hooks, Heritage Skills bridges)
\end{asciibox}

\subsubsection{Minecraft Simulation Scaling}

The PNW world is scaled so that the full vertical range of the region --- from the deepest point of Puget Sound to the summit of Mount Rainier --- maps exactly to Minecraft's build range.

\begin{center}
\rowcolors{2}{rowgray}{white}
\begin{tabularx}{\textwidth}{L{4.5cm}XC{3cm}}
\toprule
\rowcolor{primary}
\tableheader{Parameter} & \tableheader{Real World} & \tableheader{Minecraft} \\
\midrule
Mt.\ Rainier Summit & 14,410 ft (4,392 m) & y = 319 (build limit) \\
Sea Level & 0 ft & y = $-$41 \\
Puget Sound Max Depth & $-$930 ft ($-$283 m) & y = $-$64 (bedrock) \\
Total Vertical Range & 15,340 ft (4,675 m) & 383 blocks \\
Scale Factor & --- & 40.05 ft/block (12.2 m/block) \\
Horizontal Scale & TBD (approx.) & 1 block = 40 ft \\
\bottomrule
\end{tabularx}
\end{center}

\textbf{Elevation Band Mapping:}

\begin{center}
\rowcolors{2}{rowgray}{white}
\begin{tabularx}{\textwidth}{L{3.5cm}C{2.5cm}C{2.5cm}X}
\toprule
\rowcolor{secondary}
\tableheader{Life Zone} & \tableheader{Elevation (ft)} & \tableheader{Minecraft Y} & \tableheader{Biome Mapping} \\
\midrule
Arctic-Alpine & 9,500--14,410 & 196--319 & Snow, packed ice, stone \\
Subalpine (Hudsonian) & 5,000--9,500 & 84--196 & Spruce, subalpine meadow \\
Canadian (Montane) & 3,000--5,000 & 34--84 & Dense conifer forest \\
Humid Transition & 0--3,000 & $-$41--34 & Old-growth, cedar, fir \\
Intertidal & $-$15--0 & $-$41 & Sand, gravel, kelp \\
Shallow Marine & $-$200--$-$15 & $-$46--$-$41 & Eelgrass, rockfish \\
Deep Marine & $-$930--$-$200 & $-$64--$-$46 & Deep basins, dark \\
\bottomrule
\end{tabularx}
\end{center}

\subsection{Research Modules}

The taxonomy is organized into six primary research modules, each producing a comprehensive survey that feeds both the research document and the chipset derivation.

\subsubsection{Module 1: Flora (Flora-Chip)}

Vascular plants, bryophytes, and marine macroalgae of the PNW bioregion. Organized by elevation band and habitat type. Key taxa include western red cedar (\emph{Thuja plicata}), Douglas fir (\emph{Pseudotsuga menziesii}), Sitka spruce (\emph{Picea sitchensis}), western hemlock (\emph{Tsuga heterophylla}), Pacific madrone (\emph{Arbutus menziesii}), bull kelp (\emph{Nereocystis luetkeana}), eelgrass (\emph{Zostera marina}), and the alpine wildflower meadows of Rainier's subalpine zone. Cultural cross-reference: cedar is the ``Tree of Life'' to all PNW Coast peoples --- material culture, spiritual practice, and ecological identity are inseparable.

\subsubsection{Module 2: Fauna --- Terrestrial (Fauna-T-Chip)}

Mammals, birds, amphibians, and reptiles. Includes 65+ mammal species (black bear, mountain goat, Roosevelt elk, Pacific fisher, snowshoe hare, Townsend's big-eared bat), 200+ bird species (bald eagle, spotted owl, marbled murrelet, varied thrush, rufous hummingbird), 25+ amphibians (Pacific giant salamander, red-legged frog, tailed frog, rough-skinned newt --- densities up to 744 salamanders/hectare in old-growth), and 20+ reptiles. Cultural cross-reference: mountain goat wool for Salish weaving, eagle feathers in ceremony, salmon-bear nutrient pathway.

\subsubsection{Module 3: Fauna --- Aquatic \& Marine (Fauna-M-Chip)}

All five Pacific salmon species (Chinook, Coho, Sockeye, Pink, Chum), steelhead, Pacific herring, rockfish (multiple species, some endangered), lingcod, Pacific cod, spiny dogfish, plus marine mammals --- Southern Resident orcas (endangered, J/K/L pods), gray whales, humpback whales, harbor seals, Steller sea lions, sea otters --- and 3,000+ invertebrate species including Dungeness crab, geoduck clam, Pacific oyster, giant Pacific octopus, sea urchins, and sea stars. Cultural cross-reference: salmon as the economic and spiritual backbone of all PNW peoples, reef net fishing (Lummi/Saanich), whale hunting (Makah/Nuu-chah-nulth).

\subsubsection{Module 4: Fungi \& Microbiome (Myco-Chip)}

Mycorrhizal networks connecting old-growth forest root systems (the ``Wood Wide Web''), decomposer fungi recycling nurse logs, lichens as air quality indicators and ecological pioneers, soil microbiome communities, aquatic microbiome of Puget Sound. Key genera: \emph{Armillaria} (honey mushroom --- largest organisms on Earth by area), \emph{Cantharellus} (chanterelle), \emph{Boletus}, \emph{Lobaria} (lung lichen), \emph{Usnea}. The mycorrhizal network is the PNW forest's communication infrastructure --- the biological equivalent of the GSD bus architecture.

\subsubsection{Module 5: Ecological Networks (Bus-Architecture)}

Cross-cutting relationships that connect the taxonomy into a functioning ecosystem. Salmon nutrient pathway (marine $\rightarrow$ freshwater $\rightarrow$ forest via bear/eagle predation), mycorrhizal resource sharing, pollinator networks, predator-prey cascades (sea otter $\rightarrow$ sea urchin $\rightarrow$ kelp forest), watershed connectivity (headwater to estuary), fire ecology and succession, and climate-driven range shifts. Each network becomes a bus route in the derived chipset.

\subsubsection{Module 6: Cultural-Ecological Integration (Heritage-Bridge)}

The intersection of ecological knowledge and cultural practice. Ethnobotany (cedar, camas, wapato, salal, huckleberry), ethnoecology (fire management, selective harvesting, reef net engineering), traditional ecological knowledge (TEK) as a complement to Western science, and the Heritage Skills badge trail cross-references. This module bridges the PNW research block to the Salish Sea expansion of the Heritage Skills pack. All references are nation-specific per established protocol.

\subsection{Chipset Configuration}

\begin{asciibox}
# pnw-ecosystem.chipset.yaml
name: pnw-living-systems
version: 1.0.0
archetype: research-taxonomy
description: >
  PNW ecosystem research chipset derived from taxonomy.
  Each chip corresponds to a major taxonomic group.
  Bus routes correspond to ecological relationships.

math_coprocessor:
  enabled: true
  gpu: rtx-4060-ti-8gb
  engines:
    - elevation_gradient    # 15,340 ft vertical transect
    - watershed_topology    # River network graph processing
    - population_dynamics   # Lotka-Volterra, carrying capacity
    - fractal_geometry      # L-systems for branching structures
    - spatial_projection    # Geographic to Minecraft coordinates

chips:
  flora:
    agent: CRAFT-FLORA
    skills: [botany, ethnobotany, forest-ecology]
    memory_banks: [elevation-bands, habitat-types, cultural-significance]

  fauna_terrestrial:
    agent: CRAFT-FAUNA-T
    skills: [wildlife-ecology, ornithology, herpetology]
    memory_banks: [elevation-bands, habitat-types, conservation-status]

  fauna_marine:
    agent: CRAFT-FAUNA-M
    skills: [marine-biology, fisheries-science, cetology]
    memory_banks: [depth-zones, habitat-types, conservation-status]

  fungi_microbiome:
    agent: CRAFT-MYCO
    skills: [mycology, soil-science, microbial-ecology]
    memory_banks: [substrate-types, symbiotic-networks]

  ecological_networks:
    agent: CRAFT-ECO-NET
    skills: [systems-ecology, network-analysis, nutrient-cycling]
    memory_banks: [all-chips-shared]    # cross-cutting access

  heritage_bridge:
    agent: CRAFT-HERITAGE
    skills: [ethnobotany, cultural-ecology, heritage-skills]
    memory_banks: [cultural-significance, nation-specific-refs]
    safety_warden: cultural-sovereignty  # OCAP/CARE/UNDRIP

bus_architecture:
  salmon_nutrient_pathway:
    connects: [fauna_marine, fauna_terrestrial, flora, fungi_microbiome]
    direction: bidirectional
    bandwidth: high  # central ecological relationship

  mycorrhizal_network:
    connects: [flora, fungi_microbiome]
    direction: bidirectional
    bandwidth: high

  predator_prey_cascade:
    connects: [fauna_terrestrial, fauna_marine, flora]
    direction: directional
    bandwidth: medium

  watershed_connectivity:
    connects: [all]
    direction: downstream
    bandwidth: high

  cultural_ecological:
    connects: [heritage_bridge, flora, fauna_marine, fauna_terrestrial]
    direction: bidirectional
    bandwidth: medium
    safety: cultural-sovereignty-warden

clock_domains:
  alpine:       { elevation: "9500-14410", rate: slow }
  subalpine:    { elevation: "5000-9500",  rate: medium }
  montane:      { elevation: "3000-5000",  rate: medium }
  lowland:      { elevation: "0-3000",     rate: fast }
  intertidal:   { elevation: "-15-0",      rate: tidal }
  marine:       { elevation: "-930--15",   rate: current }

token_budget:
  total: 200000
  per_chip: 30000
  shared_layers: 20000
  synthesis: 20000
  minecraft_layer: 10000

minecraft_simulation:
  world_name: "Tahoma-Salish-Living-Systems"
  scale_factor_ft_per_block: 40.05
  sea_level_y: -41
  summit_y: 319
  bedrock_y: -64
  biome_mapping: elevation-band-driven
  entity_spawning: taxonomy-accurate
  redstone_systems:
    - salmon_run_indicator
    - tidal_cycle_clock
    - mycorrhizal_signal_network
\end{asciibox}

\subsection{Scope Boundaries}

\textbf{In Scope (v1.0):}
\begin{itemize}[leftmargin=2em]
\item Full taxonomy survey of PNW bioregion (western Cascades to Pacific, Yakutat to Cape Mendocino)
\item GPU math co-processor initialization for ecological computation
\item Map-reduce parallel research execution across six taxonomy modules
\item Chipset derivation from completed taxonomy with bus architecture
\item Engineering optimization pass for chipset performance
\item Minecraft world specification with geographic scaling
\item College Structure biology department seed data
\item Heritage Skills pack integration (Salish Sea expansion cross-references)
\item 59 threatened/endangered species profiles with ESA status
\item Five Pacific salmon species complete lifecycle documentation
\item Mycorrhizal network topology mapping
\end{itemize}

\textbf{Out of Scope (Future):}
\begin{itemize}[leftmargin=2em]
\item East-of-Cascades steppe/desert ecosystems (separate chipset)
\item Paleontological record (separate research mission)
\item Climate change projection modeling (requires dedicated mission with agency data)
\item Minecraft world \emph{generation} code (this mission produces the specification; a separate engineering mission builds the world generator)
\item Full Minecraft mod implementation (separate engineering pack)
\item Real-time sensor integration (future IoT bridge)
\end{itemize}

\subsection{Success Criteria}

\begin{enumerate}[leftmargin=2em]
\item The math co-processor correctly computes elevation-to-Minecraft-Y mappings for all six life zones, with scale factor verified at 40.05 ft/block.

\item Each taxonomy module produces a comprehensive species survey with: scientific name, common name, elevation range, habitat type, ecological role, conservation status, and cultural significance (where applicable).

\item Shared attribute layers (elevation bands, habitat types, ecological roles) are computed exactly once and successfully reused across all six research modules, reducing total research tokens by at least 40\% compared to independent per-species research.

\item The derived chipset definition is a valid \texttt{chipset.yaml} that passes GSD's chipset validation script and correctly routes research queries to the appropriate specialist chip.

\item Bus routes in the chipset correspond to documented ecological relationships with at least three source citations per route.

\item The Minecraft world specification produces biome boundaries at the correct Y-levels when applied to a world generator, placing sea level at y=$-$41 and the Rainier summit at y=319.

\item The College Structure biology department receives seed data organized as explorable code, with cross-references to the math department (population dynamics, fractal geometry) and heritage skills department (ethnobotany).

\item Heritage Skills integration correctly maps species to nation-specific cultural practices using nation names (never generic ``Indigenous peoples''), verified against the Salish Sea expansion vision document's geographic scope.

\item The engineering optimization pass reduces cross-chip bus contention by at least 25\% through shared memory bank reorganization.

\item All 59 ESA-listed PNW species are documented with current listing status, critical habitat designation, and primary threats.

\item The complete research block is self-contained: given only this document, an agent team can produce the full taxonomy without additional context.

\item The mission package integrates with gsd-skill-creator's observation loop, generating at least one reusable skill per taxonomy module (e.g., ``species-profile-template'', ``elevation-band-classifier'').
\end{enumerate}

\subsection{Relationship to Other Documents}

\begin{center}
\rowcolors{2}{rowgray}{white}
\begin{tabularx}{\textwidth}{L{4.5cm}X}
\toprule
\rowcolor{primary}
\tableheader{Document} & \tableheader{Relationship} \\
\midrule
gsd-silicon-layer-spec.md & Math co-processor architecture; GPU hardware specs; LoRA adapter pipeline \\
gsd-chipset-architecture-vision.md & Chipset definition format; ASIC/FPGA modes; agent topology patterns \\
gsd-skill-creator-analysis.md & Observation pipeline; pattern detection; bounded learning constraints \\
04-salish-sea-expansion-vision.md & Heritage Skills PNW expansion; nation-specific protocols; badge trail \\
cooking-with-claude-compiled-session.md & College Structure; Rosetta Core; biology department seeding \\
gsd-instruction-set-architecture-vision.md & ISA token-efficient communication; FPGA abstraction layer \\
gsd-os-desktop-vision.md & Minecraft as visualization metaphor; block-type reference \\
gsd-amiga-vision-architectural-leverage.md & The Amiga Principle; community adapter sharing \\
unit-circle-skill-creator-synthesis.md & Mathematical foundations; tangent-line activation scoring \\
The Space Between (textbook) & Fractal geometry; L-systems; complex plane of experience \\
The Quark and the Compiler & NVIDIA/CUDA history; GPU co-processor philosophy \\
\bottomrule
\end{tabularx}
\end{center}

\subsection{The Through-Line}

The salmon knows the way home. It hatches in a mountain stream fed by glacial melt from Tahoma's flanks, descends through old-growth forests connected by mycorrhizal networks, passes through estuaries where eelgrass meadows shelter juvenile fish, enters the Salish Sea where orcas and seals and eagles wait, then swims thousands of miles into the open Pacific --- and years later, returns. Returns to the exact stream. Returns carrying marine nutrients in its body that will feed the forest when the bears and eagles carry its remains onto the forest floor.

That journey --- summit to sea to ocean and back --- is the architecture of this research block. Every elevation band, every species, every ecological relationship, every cultural practice exists along that path. The math co-processor computes the geometry. The taxonomy maps the territory. The chipset routes the intelligence. The Minecraft world renders it at human scale. The college teaches it. The heritage skills honor the people who understood it first.

The spaces between the species are where the ecosystem lives. The salmon nutrient pathway is a space between marine and terrestrial. The mycorrhizal network is a space between individual trees. The potlatch is a space between competing villages. The reef net is a space between cooperating families. This research block documents those spaces --- and in documenting them, becomes one itself.

\newpage

% ============================================================
% STAGE 2 --- RESEARCH REFERENCE
% ============================================================
\stageheader{STAGE 2 --- RESEARCH REFERENCE}{secondary}

\section{PNW Living Systems --- Research Reference}

\subsection{How to Use This Document}

This reference provides the foundational data for each research module. Research agents should load the relevant module section plus the shared attribute layers before beginning species-level research. The shared layers eliminate redundant habitat and elevation research across modules.

\textbf{Key Source Organizations:}
\begin{itemize}[leftmargin=2em]
\item USFS Pacific Northwest Research Station (PNW)
\item NOAA Northwest Fisheries Science Center
\item USGS Cascades Volcano Observatory
\item Washington Department of Fish \& Wildlife (WDFW)
\item Oregon Department of Fish \& Wildlife (ODFW)
\item National Park Service (Mount Rainier, Olympic, North Cascades)
\item University of Washington School of Aquatic \& Fishery Sciences
\item Puget Sound Partnership (ecosystem recovery)
\item Burke Museum (natural history collections)
\end{itemize}

\subsection{Shared Attribute Layers}

These layers are computed once by the math co-processor and cached for all downstream modules.

\subsubsection{Elevation Bands}

\begin{center}
\rowcolors{2}{rowgray}{white}
\begin{tabularx}{\textwidth}{L{3cm}C{2cm}C{2cm}X}
\toprule
\rowcolor{secondary}
\tableheader{Zone} & \tableheader{Range (ft)} & \tableheader{MC Y-Range} & \tableheader{Dominant Community} \\
\midrule
Arctic-Alpine & 9,500+ & 196--319 & Krummholz, alpine meadow, bare rock, permanent snow/ice \\
Subalpine & 5,000--9,500 & 84--196 & Subalpine fir, mountain hemlock, avalanche lily meadows \\
Montane & 3,000--5,000 & 34--84 & Pacific silver fir, noble fir, Alaska yellow cedar \\
Lowland Forest & 500--3,000 & $-$29--34 & Douglas fir, western red cedar, western hemlock, bigleaf maple \\
Riparian/Estuary & 0--500 & $-$41--$-$29 & Red alder, Sitka spruce, willow, sedge, eelgrass transition \\
Intertidal & $-$15--0 & $-$41 & Barnacles, mussels, anemones, rockweed, sea lettuce \\
Shallow Marine & $-$200--$-$15 & $-$46--$-$41 & Eelgrass, kelp forest, rockfish habitat \\
Deep Marine & $-$930--$-$200 & $-$64--$-$46 & Soft sediment, deep basin, cold water adapted species \\
\bottomrule
\end{tabularx}
\end{center}

\subsubsection{Habitat Types (Major Categories)}

Old-Growth Conifer Forest, Second-Growth Forest, Alpine Meadow, Subalpine Parkland, Riparian Corridor, Freshwater Stream/River, Freshwater Lake, Wetland/Bog, Estuary, Rocky Intertidal, Sandy Beach, Eelgrass Meadow, Kelp Forest, Open Pelagic, Deep Basin, Volcanic/Geothermal, Urban Interface. Each habitat type has a standard profile including: typical elevation range, dominant vegetation, indicator species, primary threats, and conservation status.

\subsubsection{Ecological Roles}

Keystone Species (disproportionate ecosystem impact relative to abundance), Apex Predator, Mesopredator, Primary Producer, Primary Consumer (herbivore), Secondary Consumer, Decomposer, Pollinator, Seed Disperser, Ecosystem Engineer (beaver, salmon), Indicator Species (canary species for ecosystem health), Foundation Species (structurally defines habitat --- old-growth trees, kelp, eelgrass), Nurse Species (supports regeneration of other species --- nurse logs).

\subsection{Module 1: Flora Reference}

\subsubsection{Conifer Dominants}

The PNW is defined by its conifers. Douglas fir (\emph{Pseudotsuga menziesii}) is the most commercially important timber species in North America and the dominant tree of mid-elevation PNW forests. Western red cedar (\emph{Thuja plicata}) --- Tahoma's ``Tree of Life'' --- reaches 200+ feet and 1,000+ years; its rot-resistant wood provides the material foundation for Coast Salish culture (canoes, longhouses, clothing, baskets, ceremonial objects). Sitka spruce (\emph{Picea sitchensis}) dominates the coastal fog belt. Western hemlock (\emph{Tsuga heterophylla}) is the climax species of the lowland forest. Pacific silver fir (\emph{Abies amabilis}) and subalpine fir (\emph{Abies lasiocarpa}) define the montane and subalpine zones. Mountain hemlock (\emph{Tsuga mertensiana}) is the treeline species on Rainier.

\subsubsection{Deciduous and Understory}

Bigleaf maple (\emph{Acer macrophyllum}), red alder (\emph{Alnus rubra}) --- a nitrogen fixer critical for soil recovery after disturbance, Pacific madrone (\emph{Arbutus menziesii}) --- an evergreen broadleaf, vine maple (\emph{Acer circinatum}), Oregon grape (\emph{Mahonia aquifolium}) --- state flower of Oregon, salal (\emph{Gaultheria shallon}) --- dense understory shrub used by Coast Salish peoples for food and medicine, sword fern (\emph{Polystichum munitum}) --- the defining understory fern of PNW forests, and huckleberry species (\emph{Vaccinium} spp.) --- culturally significant food source across all PNW nations.

\subsubsection{Marine Plants}

Bull kelp (\emph{Nereocystis luetkeana}) forms canopy forests in the Salish Sea, providing critical habitat for juvenile fish, invertebrates, and marine mammals. Eelgrass (\emph{Zostera marina}) meadows in shallow bays are nursery habitat for salmon, herring, and crab. Both are declining due to water quality degradation, warming temperatures, and physical disturbance.

\subsection{Module 2: Fauna --- Terrestrial Reference}

\subsubsection{Mammals}

The PNW hosts 65+ mammal species across all elevation bands. Black bear (\emph{Ursus americanus}) --- the primary terrestrial vector for marine-derived nutrients from salmon carcasses into forests. Mountain goat (\emph{Oreamnos americanus}) --- alpine specialist, source of wool for Salish weaving. Roosevelt elk (\emph{Cervus canadensis roosevelti}) --- largest elk subspecies, Olympic Peninsula endemic population. Pacific fisher (\emph{Pekania pennanti}) --- recently reintroduced to the southern Cascades and Olympics after near-extirpation. Townsend's chipmunk, Douglas squirrel, snowshoe hare, hoary marmot (alpine), pika (\emph{Ochotona princeps}) --- a climate change indicator species whose range is retreating upward as temperatures warm.

\subsubsection{Birds}

Over 200 species documented. Bald eagle (\emph{Haliaeetus leucocephalus}) --- recovered from endangered status, a keystone scavenger on salmon carcasses. Northern spotted owl (\emph{Strix occidentalis caurina}) --- threatened, old-growth dependent, range shrinking due to barred owl competition and habitat loss. Marbled murrelet (\emph{Brachyramphus marmoratus}) --- threatened seabird that nests in old-growth trees up to 50 miles inland. Varied thrush (\emph{Ixoreus naevius}) --- the iconic song of PNW old-growth forests. Rufous hummingbird --- longest migration of any hummingbird relative to body size. Steller's jay, gray jay, American dipper (aquatic songbird), peregrine falcon (recovered).

\subsubsection{Amphibians and Reptiles}

The PNW west of the Cascades is an amphibian biodiversity hotspot with 25+ species. Plethodontid salamander densities reach 364--744 per hectare in old-growth forests (USFS PNW Research Station data). Pacific giant salamander (\emph{Dicamptodon tenebrosus}) --- one of few salamanders that vocalizes. Rough-skinned newt (\emph{Taricha granulosa}) --- produces tetrodotoxin, one of nature's most potent neurotoxins. Tailed frog (\emph{Ascaphus truei}) --- an ancient lineage adapted to cold, fast-flowing mountain streams. Oregon spotted frog --- threatened, endemic to the Pacific Northwest. Reptile diversity is lower, concentrated in oak woodland and rocky habitats.

\subsection{Module 3: Fauna --- Aquatic \& Marine Reference}

\subsubsection{Pacific Salmon (Genus \emph{Oncorhynchus})}

The five species and their characteristics:

\begin{center}
\rowcolors{2}{rowgray}{white}
\begin{tabularx}{\textwidth}{L{2.5cm}L{2.5cm}C{2cm}X}
\toprule
\rowcolor{secondary}
\tableheader{Species} & \tableheader{Common Name} & \tableheader{Max Size} & \tableheader{Cultural/Ecological Notes} \\
\midrule
\emph{O. tshawytscha} & Chinook (King) & 58 lbs & Primary prey of Southern Resident orcas; highest cultural value \\
\emph{O. kisutch} & Coho (Silver) & 12 lbs & Important sport and subsistence fish; ESA threatened \\
\emph{O. nerka} & Sockeye (Red) & 8 lbs & Lake-rearing life history; Baker Lake runs \\
\emph{O. gorbuscha} & Pink (Humpback) & 5 lbs & Two-year cycle; largest total biomass \\
\emph{O. keta} & Chum (Dog) & 15 lbs & Latest spawner; important for smoking/drying \\
\emph{O. mykiss} & Steelhead & 20 lbs & Anadromous rainbow trout; ESA listed \\
\bottomrule
\end{tabularx}
\end{center}

\subsubsection{Marine Mammals}

Southern Resident Killer Whales (J, K, L pods): approximately 74 individuals as of recent surveys, ESA endangered, primary threats are Chinook salmon decline, vessel noise, and contaminant accumulation. Gray whales migrate through annually. Humpback whales increasingly present in the Salish Sea. Harbor seals are the most common marine mammal. Steller sea lions (ESA threatened in western DPS). Sea otters --- historically extirpated from Washington, partial recovery underway on the Olympic coast.

\subsubsection{Marine Invertebrates}

Over 3,000 invertebrate species in Puget Sound. Dungeness crab (\emph{Metacarcinus magister}) --- the commercially and culturally most important crustacean. Geoduck clam (\emph{Panopea generosa}) --- the largest burrowing clam in the world, living 140+ years. Giant Pacific octopus (\emph{Enteroctopus dofleini}) --- the world's largest octopus species, found throughout the Sound. Sunflower sea star (\emph{Pycnopodia helianthoides}) --- critically declining due to sea star wasting syndrome, a keystone predator controlling sea urchin populations and thus kelp forest health.

\subsection{Module 4: Fungi \& Microbiome Reference}

Old-growth PNW forests contain some of the most extensive mycorrhizal networks documented. Individual \emph{Armillaria} (honey mushroom) clones can span thousands of acres. Ectomycorrhizal fungi connect the root systems of conifers into resource-sharing networks --- the ``Wood Wide Web'' --- through which trees share carbon, water, and chemical warning signals. Mother trees (large, old individuals) serve as network hubs. Douglas fir seedlings connected to mother trees via mycorrhizal networks have significantly higher survival rates than isolated seedlings.

Chanterelle mushrooms (\emph{Cantharellus formosus} --- the Pacific golden chanterelle) are commercially harvested throughout the PNW and are ectomycorrhizal partners with Douglas fir and western hemlock. The PNW matsutake (\emph{Tricholoma murrillianum}) has deep cultural significance for Japanese-American communities and is an important non-timber forest product.

\subsection{Module 5: Ecological Networks Reference}

\subsubsection{Salmon Nutrient Pathway}

Marine-derived nitrogen (MDN) transported by salmon into freshwater and terrestrial ecosystems is one of the most well-documented nutrient pathways in ecology. Bears, eagles, and other predators carry salmon carcasses onto the forest floor, where decomposition releases nitrogen-15 (a marine isotope marker) into soil. Studies using nitrogen isotope analysis have found marine-derived nitrogen in trees, shrubs, insects, and songbirds up to 500 meters from salmon-bearing streams. This pathway connects the deep Pacific Ocean to the alpine forests of the Cascades through a chain of biological carriers.

\subsubsection{Kelp-Otter-Urchin Cascade}

A classic trophic cascade: sea otters control sea urchin populations; without otters, urchins overgraze kelp, creating ``urchin barrens.'' Where sea otters have been restored, kelp forests recover, providing habitat for hundreds of species. The loss of sunflower sea stars to wasting syndrome has created a parallel urchin control crisis.

\subsubsection{Fire Ecology}

Low-intensity fire historically maintained open prairies and oak woodlands in the Puget Sound lowlands and Willamette Valley. Many Coast Salish and Kalapuya nations used controlled burning to manage camas prairies, encourage berry production, and maintain hunting grounds. Fire suppression since European settlement has converted these landscapes to closed-canopy forest, reducing biodiversity. Prescribed fire is now recognized as essential for ecosystem health.

\subsection{Module 6: Cultural-Ecological Integration Reference}

This module bridges to the Heritage Skills Salish Sea expansion. Research must follow established protocols: nation-specific references only (never generic ``Indigenous peoples''), OCAP/CARE/UNDRIP cultural sovereignty framework, and community-authorized sources where available. Key ethnobotanical relationships include: western red cedar (Coast Salish --- material culture foundation), camas (\emph{Camassia quamash}) --- a dietary staple managed through controlled burning, wapato (\emph{Sagittaria latifolia}) --- an aquatic tuber harvested from wetlands, and salal berries --- mixed with eulachon grease as a preserved food. These relationships are documented in the Heritage Skills badge trail under the Cedar Path and Salmon Path.

\subsection{Safety and Sensitivity Considerations}

\begin{center}
\rowcolors{2}{rowgray}{white}
\begin{tabularx}{\textwidth}{L{3.5cm}C{2cm}X}
\toprule
\rowcolor{tertiary!80!black}
\tableheader{Consideration} & \tableheader{Boundary} & \tableheader{Protocol} \\
\midrule
Endangered species locations & ABSOLUTE & Never publish precise GPS coordinates of nest sites, dens, or spawning locations for ESA-listed species \\
Cultural sovereignty & ABSOLUTE & All Indigenous ecological knowledge attributed to specific nations; OCAP/CARE/UNDRIP framework applies \\
Potlatch sensitivity & GATE & Note historical criminalization (1884--1951 in Canada); present with appropriate gravity \\
TEK integration & GATE & Present as complementary knowledge system, not subordinate to Western science \\
Marine mammal approach & ANNOTATE & Note legal approach distances (MMPA: 100 yards for whales, 200 yards for SRKWs) \\
Harvest data & ANNOTATE & Distinguish tribal treaty rights from recreational harvest; do not present as equivalent \\
\bottomrule
\end{tabularx}
\end{center}

\subsection{Source Bibliography}

\subsubsection{Government \& Agency Sources}

USFS Pacific Northwest Research Station --- wildlife ecology, amphibian surveys, old-growth forest research. NOAA Northwest Fisheries Science Center --- salmon, marine mammals, ecosystem analysis. NOAA Fisheries West Coast Region --- ESA listings, critical habitat designations. National Park Service --- Mount Rainier, Olympic, North Cascades natural resource inventories. USGS --- volcanic monitoring, water resources, species accounts. Washington Department of Fish \& Wildlife --- state species of concern, harvest data. Puget Sound Partnership --- ecosystem indicators, recovery targets. U.S. Fish \& Wildlife Service --- ESA species profiles, recovery plans.

\subsubsection{Peer-Reviewed Research}

Bury \& Corn (1987, 1988) --- amphibian densities in Douglas-fir forests. Helfield \& Naiman (2001) --- marine-derived nitrogen in riparian forests. Simard et al.\ (1997, 2012) --- mycorrhizal networks and resource sharing in forests. Quinn (2005) --- \emph{The Behavior and Ecology of Pacific Salmon and Trout}. Franklin \& Dyrness (1973, revised 1988) --- \emph{Natural Vegetation of Oregon and Washington}. DellaSala (2011) --- temperate rainforest biodiversity and ecosystem services.

\subsubsection{Cultural \& Ethnographic Sources}

Stewart, Hilary --- \emph{Cedar: Tree of Life} (1984); \emph{Indian Fishing} (1977). Suttles, Wayne --- \emph{Coast Salish Essays} (1987); \emph{Handbook of North American Indians Vol.\ 7} (1990). Turner, Nancy J.\ --- \emph{Ancient Pathways, Ancestral Knowledge} (2 vols., 2014). See Heritage Skills pack bibliography for complete ethnographic references.

\newpage

% ============================================================
% STAGE 3 --- MISSION PACKAGE
% ============================================================
\stageheader{STAGE 3 --- MISSION PACKAGE}{tertiary}

\section{Milestone Specification}

\metafield{Mission Objective:}{Produce a comprehensive PNW ecosystem taxonomy organized through GPU-accelerated computation, derived into a GSD chipset architecture, integrated with the College Structure biology department and Heritage Skills educational pack, and specified as a geographically accurate Minecraft simulation world. The output is both a research reference and a functioning chipset definition ready for GSD orchestrator execution.}

\subsection{Architecture Overview}

\begin{asciibox}
                    MISSION ARCHITECTURE
                    ====================

  Wave 0: Foundation           Wave 1: Parallel Research
  (Sequential)                 (3 Parallel Tracks)
  +-----------------+         +-------------------+
  | Shared Types    |-------->| Track A: Flora +  |
  | Attribute Layers|    |    |   Fungi/Microbiome|
  | Math Co-Proc    |    |    +-------------------+
  | Species Template|    |    +-------------------+
  +-----------------+    +--->| Track B: Fauna    |
                         |    |   Terrestrial +   |
                         |    |   Marine/Aquatic   |
                         |    +-------------------+
                         |    +-------------------+
                         +--->| Track C: Eco Nets |
                              |   + Heritage      |
                              |   Bridge          |
                              +-------------------+
                                       |
                              Wave 2: Synthesis (Sequential)
                              +-------------------+
                              | Cross-module merge|
                              | Chipset derivation|
                              | Bus architecture  |
                              +-------------------+
                                       |
                              Wave 3: Engineering (Sequential)
                              +-------------------+
                              | Optimization pass |
                              | Minecraft spec    |
                              | College integratn |
                              | Verification      |
                              +-------------------+
                                       |
                              Wave 4: Publication (Sequential)
                              +-------------------+
                              | Final assembly    |
                              | Test execution    |
                              | Package delivery  |
                              +-------------------+
\end{asciibox}

\subsection{System Layers}

\begin{enumerate}[leftmargin=2em]
\item \textbf{Math Co-Processor Layer} --- GPU-accelerated computation for elevation gradients, watershed topology, population dynamics, fractal geometry, and Minecraft coordinate projection.
\item \textbf{Shared Attribute Layer} --- Cached elevation bands, habitat types, ecological roles, conservation status, and cultural significance profiles reused across all taxonomy modules.
\item \textbf{Taxonomy Research Layer} --- Six parallel research modules (Flora, Fauna-T, Fauna-M, Fungi, Eco-Nets, Heritage Bridge) producing species profiles against shared attributes.
\item \textbf{Chipset Derivation Layer} --- Transforms completed taxonomy into chipset.yaml with chips, buses, memory banks, and clock domains.
\item \textbf{Engineering Optimization Layer} --- Performance tuning of derived chipset for cache coherence, bus contention, and token efficiency.
\item \textbf{Simulation Layer} --- Minecraft world specification with elevation-accurate biome mapping.
\item \textbf{Integration Layer} --- College Structure biology department seeding and Heritage Skills cross-references.
\end{enumerate}

\subsection{Deliverables}

\begin{center}
\rowcolors{2}{rowgray}{white}
\begin{tabularx}{\textwidth}{C{0.5cm}L{3.5cm}X}
\toprule
\rowcolor{primary}
\tableheader{\#} & \tableheader{Deliverable} & \tableheader{Acceptance Criteria} \\
\midrule
1 & Math Co-Processor Config & silicon.yaml with 5 engine definitions; validated against RTX 4060 Ti specs \\
2 & Shared Attribute Layers & 8 elevation bands, 17 habitat types, 13 ecological roles --- all with Minecraft Y mapping \\
3 & Flora Survey & 50+ species profiles with shared-attribute cross-references \\
4 & Fauna-Terrestrial Survey & 80+ species (mammals, birds, amphibians, reptiles) \\
5 & Fauna-Marine Survey & 60+ species (salmon, marine mammals, invertebrates) \\
6 & Fungi/Microbiome Survey & 30+ species/communities with network topology \\
7 & Ecological Networks Map & 5 major pathways documented with citations \\
8 & Heritage Bridge Document & Nation-specific cross-references for all cultural-ecological entries \\
9 & Derived Chipset (chipset.yaml) & Valid GSD chipset with 6 chips, 5 bus routes, 6 clock domains \\
10 & Engineering Optimization Report & Documented improvements with before/after metrics \\
11 & Minecraft World Spec & Complete biome-to-elevation mapping, entity spawn tables, redstone systems \\
12 & College Biology Dept Seed & Explorable code structure with math/heritage cross-references \\
\bottomrule
\end{tabularx}
\end{center}

\subsection{Component Breakdown}

\begin{center}
\rowcolors{2}{rowgray}{white}
\begin{tabularx}{\textwidth}{L{3cm}L{2cm}C{1.5cm}C{1.5cm}}
\toprule
\rowcolor{primary}
\tableheader{Component} & \tableheader{Dependencies} & \tableheader{Model} & \tableheader{Est. Tokens} \\
\midrule
Math Co-Processor Init & None & Sonnet & 15K \\
Shared Attribute Layers & Co-Proc & Haiku & 10K \\
Flora Survey & Attributes & Sonnet & 35K \\
Fauna-T Survey & Attributes & Sonnet & 40K \\
Fauna-M Survey & Attributes & Sonnet & 35K \\
Fungi Survey & Attributes & Sonnet & 25K \\
Eco Networks & All Surveys & Opus & 30K \\
Heritage Bridge & Eco Nets + Surveys & Opus & 25K \\
Chipset Derivation & All Above & Opus & 20K \\
Engineering Optimization & Chipset & Opus & 15K \\
Minecraft World Spec & Optimized Chipset & Sonnet & 20K \\
College Integration & All Above & Sonnet & 15K \\
Verification \& Test & All Above & Sonnet & 15K \\
Final Assembly & All Above & Sonnet & 10K \\
\midrule
\multicolumn{3}{l}{\textbf{Total}} & \textbf{310K} \\
\bottomrule
\end{tabularx}
\end{center}

\subsection{Model Assignment Rationale}

\begin{center}
\rowcolors{2}{rowgray}{white}
\begin{tabularx}{\textwidth}{C{2cm}C{1.5cm}X}
\toprule
\rowcolor{secondary}
\tableheader{Model} & \tableheader{Allocation} & \tableheader{Rationale} \\
\midrule
Opus & 29\% & Ecological network synthesis, heritage bridge (cultural sensitivity), chipset derivation (architectural judgment), engineering optimization \\
Sonnet & 62\% & Species surveys (structured research with template), math co-processor config, Minecraft spec, college integration, verification \\
Haiku & 9\% & Shared attribute scaffolding, template generation, directory structure \\
\bottomrule
\end{tabularx}
\end{center}

\section{Wave Execution Plan}

\subsection{Wave Summary}

\begin{center}
\rowcolors{2}{rowgray}{white}
\begin{tabularx}{\textwidth}{L{2cm}C{1.5cm}C{2cm}C{2cm}X}
\toprule
\rowcolor{primary}
\tableheader{Wave} & \tableheader{Tasks} & \tableheader{Parallel Tracks} & \tableheader{Est. Time} & \tableheader{Cache Dependencies} \\
\midrule
Wave 0 & 2 & 1 (sequential) & 1 session & None \\
Wave 1 & 6 & 3 & 3 sessions & Wave 0 outputs \\
Wave 2 & 3 & 1 (sequential) & 2 sessions & Wave 1 all tracks \\
Wave 3 & 3 & 1 (sequential) & 2 sessions & Wave 2 outputs \\
Wave 4 & 2 & 1 (sequential) & 1 session & Wave 3 outputs \\
\midrule
\multicolumn{2}{l}{\textbf{Total: 16 tasks}} & & \textbf{9 sessions} & \\
\bottomrule
\end{tabularx}
\end{center}

\subsection{Wave 0: Foundation (Sequential)}

\begin{center}
\rowcolors{2}{rowgray}{white}
\begin{tabularx}{\textwidth}{L{3.5cm}XC{1.5cm}C{1.5cm}}
\toprule
\rowcolor{secondary}
\tableheader{Task} & \tableheader{Produces} & \tableheader{Model} & \tableheader{Est. Tokens} \\
\midrule
Math Co-Processor Init & silicon.yaml engines, coordinate projection tables & Sonnet & 15K \\
Shared Attribute Layers & Elevation bands, habitat types, ecological roles, species template & Haiku & 10K \\
\bottomrule
\end{tabularx}
\end{center}

\subsection{Wave 1: Taxonomy Map (3 Parallel Tracks)}

\textbf{Track A: Producers (Flora + Fungi)}

\begin{center}
\rowcolors{2}{rowgray}{white}
\begin{tabularx}{\textwidth}{L{2.5cm}XC{1.5cm}C{1.5cm}}
\toprule
\rowcolor{accent!80!black}
\tableheader{Task} & \tableheader{Produces} & \tableheader{Model} & \tableheader{Tokens} \\
\midrule
Flora Survey & 50+ species profiles, conifer/deciduous/marine plant sections & Sonnet & 35K \\
Fungi Survey & 30+ species/community profiles, mycorrhizal network map & Sonnet & 25K \\
\bottomrule
\end{tabularx}
\end{center}

\textbf{Track B: Consumers (Fauna Terrestrial + Marine)}

\begin{center}
\rowcolors{2}{rowgray}{white}
\begin{tabularx}{\textwidth}{L{2.5cm}XC{1.5cm}C{1.5cm}}
\toprule
\rowcolor{accent!80!black}
\tableheader{Task} & \tableheader{Produces} & \tableheader{Model} & \tableheader{Tokens} \\
\midrule
Fauna-T Survey & 80+ terrestrial species profiles across 4 classes & Sonnet & 40K \\
Fauna-M Survey & 60+ aquatic/marine species profiles & Sonnet & 35K \\
\bottomrule
\end{tabularx}
\end{center}

\textbf{Track C: Integrators (Eco Networks + Heritage Bridge)}

\begin{center}
\rowcolors{2}{rowgray}{white}
\begin{tabularx}{\textwidth}{L{2.5cm}XC{1.5cm}C{1.5cm}}
\toprule
\rowcolor{accent!80!black}
\tableheader{Task} & \tableheader{Produces} & \tableheader{Model} & \tableheader{Tokens} \\
\midrule
Ecological Networks & 5 major pathway documents with cross-module references & Opus & 30K \\
Heritage Bridge & Nation-specific cultural-ecological cross-reference table & Opus & 25K \\
\bottomrule
\end{tabularx}
\end{center}

\subsection{Wave 2: Synthesis \& Chipset Derivation (Sequential)}

\begin{center}
\rowcolors{2}{rowgray}{white}
\begin{tabularx}{\textwidth}{L{3cm}XC{1.5cm}C{1.5cm}}
\toprule
\rowcolor{secondary}
\tableheader{Task} & \tableheader{Produces} & \tableheader{Model} & \tableheader{Tokens} \\
\midrule
Cross-Module Merge & Unified taxonomy with all cross-references resolved & Opus & 10K \\
Chipset Derivation & pnw-ecosystem.chipset.yaml with 6 chips, 5 buses & Opus & 20K \\
Minecraft Mapping & Complete biome-to-Y mapping, entity spawn tables & Sonnet & 20K \\
\bottomrule
\end{tabularx}
\end{center}

\subsection{Wave 3: Engineering \& Integration (Sequential)}

\begin{center}
\rowcolors{2}{rowgray}{white}
\begin{tabularx}{\textwidth}{L{3cm}XC{1.5cm}C{1.5cm}}
\toprule
\rowcolor{secondary}
\tableheader{Task} & \tableheader{Produces} & \tableheader{Model} & \tableheader{Tokens} \\
\midrule
Engineering Optimization & Optimized chipset with performance report & Opus & 15K \\
College Integration & Biology dept seed data with cross-references & Sonnet & 15K \\
Verification \& Test & Test execution results, verification matrix & Sonnet & 15K \\
\bottomrule
\end{tabularx}
\end{center}

\subsection{Wave 4: Publication (Sequential)}

\begin{center}
\rowcolors{2}{rowgray}{white}
\begin{tabularx}{\textwidth}{L{3cm}XC{1.5cm}C{1.5cm}}
\toprule
\rowcolor{secondary}
\tableheader{Task} & \tableheader{Produces} & \tableheader{Model} & \tableheader{Tokens} \\
\midrule
Final Assembly & Complete mission package document & Sonnet & 5K \\
Package Delivery & PDF + .tex + index.html + chipset.yaml & Sonnet & 5K \\
\bottomrule
\end{tabularx}
\end{center}

\subsection{Cache Optimization Strategy}

\begin{itemize}[leftmargin=2em]
\item \textbf{Shared attribute precomputation:} Elevation bands, habitat types, and ecological roles computed in Wave 0 and cached for all Wave 1 agents. Estimated token savings: 40\% reduction in per-species habitat descriptions.
\item \textbf{Species template reuse:} A single species profile template (scientific name, common name, elevation range, habitat, role, status, cultural significance) loaded once per track.
\item \textbf{Cross-track deduplication:} Species appearing in multiple modules (e.g., salmon in both Fauna-M and Eco Networks) are profiled once in the primary module and cross-referenced elsewhere.
\item \textbf{Knowledge tiering:} Summary tier ($\sim$5K) loaded by all agents; active tier ($\sim$15K) loaded by relevant track; reference tier ($\sim$30K+) loaded only for deep dives.
\end{itemize}

\subsection{Dependency Graph}

\begin{asciibox}
  Wave 0                Wave 1                    Wave 2         Wave 3        Wave 4
  ======                ======                    ======         ======        ======

  [Co-Proc]---+------>[Flora]------+
              |                    |
              +------>[Fungi]------+------>[Merge]--+
              |                    |                 |
              +------>[Fauna-T]----+   [Chipset]<---+-->[Optimize]--+
              |                    |       ^                        |
              +------>[Fauna-M]----+       |         [College]<----+
              |                    |       |                       |
              +------>[Eco-Nets]---+-------+         [Verify]<----+-->[Publish]
              |                    |                               |
  [Attribs]---+------>[Heritage]---+------>[MC-Map]--+            |
                                                     +-->----------+

  CRITICAL PATH: Co-Proc -> Flora/Fauna-T(parallel) -> Merge -> Chipset
                 -> Optimize -> Verify -> Publish
\end{asciibox}

\subsection{Token Budget}

\begin{center}
\rowcolors{2}{rowgray}{white}
\begin{tabularx}{\textwidth}{L{3cm}C{2cm}C{2.5cm}C{2cm}}
\toprule
\rowcolor{primary}
\tableheader{Wave} & \tableheader{Est. Tokens} & \tableheader{Context Windows} & \tableheader{Sessions} \\
\midrule
Wave 0 & 25K & 2 & 1 \\
Wave 1 & 190K & 12 & 4 \\
Wave 2 & 50K & 4 & 2 \\
Wave 3 & 45K & 4 & 2 \\
Wave 4 & 10K & 1 & 1 \\
\midrule
\textbf{Total} & \textbf{320K} & \textbf{23} & \textbf{10} \\
\bottomrule
\end{tabularx}
\end{center}

\section{Test Plan}

\subsection{Test Categories}

\begin{center}
\rowcolors{2}{rowgray}{white}
\begin{tabularx}{\textwidth}{L{3.5cm}C{1.5cm}C{2cm}C{2cm}}
\toprule
\rowcolor{primary}
\tableheader{Category} & \tableheader{Count} & \tableheader{Priority} & \tableheader{Failure Action} \\
\midrule
Safety-Critical & 6 & Mandatory & BLOCK \\
Core Functionality & 12 & Required & BLOCK \\
Integration & 8 & Required & BLOCK \\
Edge Cases & 6 & Best-effort & LOG \\
\midrule
\textbf{Total} & \textbf{32} & & \\
\bottomrule
\end{tabularx}
\end{center}

\subsection{Safety-Critical Tests (Mandatory Pass)}

\begin{center}
\rowcolors{2}{rowgray}{white}
\begin{tabularx}{\textwidth}{C{1cm}L{3cm}X}
\toprule
\rowcolor{tertiary!80!black}
\tableheader{ID} & \tableheader{Verifies} & \tableheader{Expected Behavior} \\
\midrule
SC-1 & No precise ESA species locations & No GPS coordinates for endangered species nest/den/spawn sites appear anywhere in output \\
SC-2 & Nation-specific attribution & Every cultural reference uses specific nation name; zero instances of generic ``Indigenous peoples'' \\
SC-3 & Cultural sovereignty protocol & OCAP/CARE/UNDRIP framework applied to all Heritage Bridge content \\
SC-4 & Potlatch sensitivity & Criminalization history (1884--1951) noted wherever potlatch is discussed \\
SC-5 & Marine mammal distances & Legal approach distances cited for all marine mammal species \\
SC-6 & Treaty rights distinction & Tribal treaty harvest rights distinguished from recreational harvest in all fisheries references \\
\bottomrule
\end{tabularx}
\end{center}

\subsection{Verification Matrix}

\begin{center}
\rowcolors{2}{rowgray}{white}
\begin{tabularx}{\textwidth}{C{0.8cm}X L{2.5cm}}
\toprule
\rowcolor{primary}
\tableheader{SC\#} & \tableheader{Success Criterion} & \tableheader{Test ID(s)} \\
\midrule
1 & Co-proc elevation mapping accuracy & CF-1, CF-2 \\
2 & Comprehensive species surveys & CF-3, CF-4, CF-5, CF-6 \\
3 & Shared attribute reuse (40\% savings) & CF-7 \\
4 & Valid chipset.yaml & CF-8, INT-1 \\
5 & Ecological bus routes cited & CF-9 \\
6 & Minecraft biome boundaries correct & CF-10, INT-2 \\
7 & College seed data structured & INT-3 \\
8 & Heritage Skills cross-references & SC-2, INT-4 \\
9 & Engineering optimization (25\% reduction) & CF-11 \\
10 & All 59 ESA species documented & CF-12 \\
11 & Self-contained research block & INT-5 \\
12 & Reusable skills generated & INT-6 \\
\bottomrule
\end{tabularx}
\end{center}

\section{Mission Crew Manifest}

\textbf{Activation Profile: Fleet}

\begin{center}
\rowcolors{2}{rowgray}{white}
\begin{tabularx}{\textwidth}{L{2cm}L{3cm}X}
\toprule
\rowcolor{primary}
\tableheader{Role} & \tableheader{Agent} & \tableheader{Responsibility} \\
\midrule
FLIGHT & Orchestrator & Mission direction; go/no-go gates at wave boundaries \\
PLAN & Planner & Wave decomposition; parallel track optimization; cache timing \\
SCOUT & Research Agent & Source gathering from USFS, NOAA, USGS, NPS, WDFW \\
EXEC-A & Executor (Track A) & Flora + Fungi survey production \\
EXEC-B & Executor (Track B) & Fauna-T + Fauna-M survey production \\
EXEC-C & Executor (Track C) & Eco Networks + Heritage Bridge production \\
CRAFT-FLORA & Flora Specialist & Botanical taxonomy, ethnobotany \\
CRAFT-FAUNA & Fauna Specialist & Wildlife ecology, marine biology, herpetology \\
CRAFT-MYCO & Mycology Specialist & Mycorrhizal networks, soil ecology \\
CRAFT-ECO & Ecology Specialist & Systems ecology, network analysis, nutrient cycling \\
CRAFT-HERITAGE & Heritage Specialist & Cultural sovereignty, nation-specific protocols \\
VERIFY & Verification & Source validation; taxonomy accuracy checking \\
INTEG & Integration Agent & Cross-module relationships; chipset derivation \\
CAPCOM & HITL Interface & Human approval at wave boundaries \\
BUDGET & Resource Manager & Token budget tracking; session planning \\
SURGEON & Health Monitor & Research quality degradation detection \\
TOPO & Topology Manager & Agent team structure; mid-mission restructuring \\
ARCH & Architecture & Chipset structure decisions; bus design \\
SECURE & Safety Warden & Cultural sovereignty; ESA compliance; safety boundaries \\
LOG & Chronicle & Audit trail; commit messages; milestone summaries \\
\bottomrule
\end{tabularx}
\end{center}

\section{Execution Summary}

\begin{center}
\rowcolors{2}{rowgray}{white}
\begin{tabularx}{\textwidth}{L{5cm}X}
\toprule
\rowcolor{primary}
\tableheader{Metric} & \tableheader{Value} \\
\midrule
Total Tasks & 16 \\
Parallel Tracks & 3 (Wave 1) \\
Sequential Depth & 5 waves \\
Opus / Sonnet / Haiku Split & 29\% / 62\% / 9\% \\
Estimated Context Windows & 23 \\
Estimated Sessions & 10 \\
Estimated Wall Time & 8--12 hours \\
Total Tests & 32 (6 safety-critical) \\
Deliverables & 12 \\
Activation Profile & Fleet (20 roles) \\
Species Coverage Target & 220+ profiles \\
Minecraft Scale Factor & 40.05 ft/block (12.2 m/block) \\
\bottomrule
\end{tabularx}
\end{center}

\vspace{1cm}

\begin{quotebox}
\textit{The salmon knows the way home. From the glacial streams of Tahoma through the mycorrhizal forests to the deep basins of the Salish Sea, every species, every relationship, every elevation band exists along that path. The math co-processor computes the geometry. The taxonomy maps the territory. The chipset routes the intelligence. The Minecraft world renders it at human scale. The college teaches it. The heritage skills honor the people who understood it first.}

\vspace{0.5em}

\textit{The spaces between the species are where the ecosystem lives. This research block documents those spaces --- and in documenting them, becomes one itself.}

\vspace{0.5em}

\hfill --- \textit{The Through-Line}
\end{quotebox}

\end{document}
